% ===== PRÉAMBULE CANONIQUE — à coller VERBATIM en tête de CHAQUE .tex =====
\documentclass[11pt,a4paper]{article}
\usepackage[utf8]{inputenc}\usepackage[T1]{fontenc}\usepackage{lmodern}
\usepackage[french]{babel}
\usepackage{amsmath,amssymb,mathtools}\usepackage{icomma}
\usepackage{siunitx}\sisetup{locale=FR}  % \qty{}{}, \si{}, \ang{}
\usepackage{xcolor}\usepackage{colortbl}
\usepackage{tikz}\usepackage{pgfplots}\pgfplotsset{compat=1.18}
\usepackage[siunitx,europeanresistors]{circuitikz} % circuits (élec.)
\usepackage[most]{tcolorbox}\usepackage{enumitem}\usepackage{array,booktabs}
\usepackage[a4paper,margin=2.2cm]{geometry}
\usetikzlibrary{arrows.meta,calc,angles,quotes,positioning,patterns,%
  backgrounds,shapes.geometric,decorations.pathmorphing,decorations.markings,intersections}
\definecolor{primary}{RGB}{222,49,99}\definecolor{primarydark}{RGB}{150,22,55}
\definecolor{defcol}{RGB}{0,114,178}\definecolor{methcol}{RGB}{52,92,148}
\definecolor{excol}{RGB}{0,135,90}\definecolor{attncol}{RGB}{200,40,55}
\tcbset{boxbase/.style={enhanced,breakable,boxrule=0.9pt,arc=2.5pt,
  left=3mm,right=3mm,top=2.3mm,bottom=2.3mm,fonttitle=\bfseries,coltitle=white}}
% --- boîtes TUTORIEL (noms EXACTS, ne pas renommer) ---
\newtcolorbox{cle}[1][]{boxbase,colback=defcol!6,colframe=defcol,
  title={\textsc{À retenir}\ifx\relax#1\relax\relax\else\textmd{ : \itshape #1}\fi}}
\newtcolorbox{methode}[1][]{boxbase,colback=methcol!7,colframe=methcol,sharp corners,
  title={\textsc{Méthode}\ifx\relax#1\relax\relax\else\textmd{ --- \itshape #1}\fi}}
\newtcolorbox{exemplebox}[1][]{boxbase,colback=excol!5,colframe=excol,
  title={\textsc{Exemple}\ifx\relax#1\relax\relax\else\textmd{ : \itshape #1}\fi}}
\newtcolorbox{piege}[1][]{boxbase,colback=attncol!6,colframe=attncol,
  title={\textsc{Piège}\ifx\relax#1\relax\relax\else\textmd{ : \itshape #1}\fi}}
\newtcolorbox{remarque}[1][]{boxbase,colback=black!4,colframe=black!45,
  title={\textsc{Remarque}\ifx\relax#1\relax\relax\else\textmd{ : \itshape #1}\fi}}
% --- boîtes EXERCICE / CORRIGÉ (noms EXACTS, auto-comptées) ---
\newtcolorbox[auto counter]{exo}[1][]{boxbase,colback=primary!4,colframe=primary,
  title={\textsc{Exercice}~\thetcbcounter\ifx\relax#1\relax\relax\else\textmd{ --- \itshape #1}\fi}}
\newtcolorbox[auto counter]{corrige}[1][]{boxbase,colback=excol!4,colframe=excol,
  title={\textsc{Corrigé}~\thetcbcounter\ifx\relax#1\relax\relax\else\textmd{ --- \itshape #1}\fi}}
\newcommand{\vv}[1]{\vec{#1}}\newcommand{\ds}{\displaystyle}
\newcommand{\R}{\mathbb{R}}\newcommand{\N}{\mathbb{N}}\newcommand{\Z}{\mathbb{Z}}
\begin{document}
\sisetup{per-mode=symbol}

\begin{exo}[vrai ou faux sur le microscope]
Réponds par \textbf{vrai} ou \textbf{faux} et justifie brièvement.
\begin{enumerate}[label=\alph*)]
  \item Pour obtenir une image à l'infini, l'image donnée par l'objectif doit se trouver dans le plan
        focal de l'oculaire.
  \item Si on remplace l'oculaire par un autre de focale \emph{double}, le grossissement est divisé
        par deux.
  \item L'image donnée par l'objectif est droite et agrandie.
  \item L'oculaire est un système divergent.
  \item L'image finale donnée par l'oculaire est virtuelle, droite et agrandie.
\end{enumerate}
\end{exo}

\begin{exo}[la loupe brisée]
Une loupe (lentille convergente) tombe et se casse en deux morceaux, un grand et un petit. Peut-on
encore s'en servir pour observer un objet ? Justifie en te rappelant que \emph{chaque} portion de
lentille reçoit de la lumière issue de tout l'objet.
\end{exo}

\begin{exo}[choisir le bon oculaire]
Un microscope a un objectif de grandissement $|\gamma_1|=30$ et un oculaire de grossissement $\times 10$.
\begin{enumerate}[label=\alph*)]
  \item Quel est le grossissement total ?
  \item On remplace l'oculaire par un autre de \textbf{focale double}. Que devient le grossissement
        total ?
  \item Pourquoi préfère-t-on, pour grossir fortement, un oculaire de \emph{courte} focale ?
\end{enumerate}
\end{exo}

\begin{exo}[photo ou projecteur ?]
Deux dispositifs emploient un objectif convergent de focale $f=\qty{5}{\cm}$.
\begin{enumerate}[label=\alph*)]
  \item Dispositif 1 : objet à $u=\qty{200}{\cm}$. Calcule $v$ et $\gamma$, décris l'image et nomme
        l'instrument.
  \item Dispositif 2 : objet à $u=\qty{6}{\cm}$. Calcule $v$ et $\gamma$, décris l'image et nomme
        l'instrument.
\end{enumerate}
\end{exo}

\end{document}
